\documentclass[11pt]{article}

\usepackage{amsmath,amssymb,amsthm,mathtools}
\usepackage[margin=1in]{geometry}
\usepackage{url}

\newtheorem{theorem}{Theorem}[section]
\newtheorem{lemma}[theorem]{Lemma}
\newcommand{\Gap}{\Gamma_H}

\begin{document}

\section{Purpose}

This note treats a singular high-density deformation outside the uniform
bounded-complement theorem.  Let $S$ be a measurable
set of mass $s$ and put
\[
 W_s(x,y)=1-\mathbf 1_S(x)\mathbf 1_S(y).
 \tag{1.1}\label{eq:hole}
\]
Thus one deletes all edges inside one concentrating set.  Its edge density is
$p_s=1-s^2$, and its normalized complement is
\[
 U_s=\frac{1-W_s}{1-p_s}=s^{-2}\mathbf 1_{S\times S}.
\]
Consequently $\lVert U_s\rVert_\infty=s^{-2}=(1-p_s)^{-1}$ diverges at
exactly the scale excluded by the bounded-complement theorem.

For all $29$ currently open Atlas rows, the exact comparison with the
chromatic target is nonnegative on the whole required density interval, and
is strict except at $p=1$.  Thus the most basic concentrating complement---a
single deleted clique---does not supply a negative construction.  This is a
one-parameter structural calculation, not an optimization or template
search.  A later theorem extends the high-density conclusion uniformly to
every finite or countable complete-multipartite partition, including the
critical number of parts and arbitrary part imbalances.  A second exact
 full-interval theorem treats every finite or countable union of clique
 blocks with arbitrary dust, including the canonical nonregular complete-cut
 and persistent-hub boundary examples.  The independent note
 \path{notes/two_block_hub_high_density_stability.tex} further fills the
 entire constant two-block hub cube of exceptional mass at most $1/3$; the
 further note
 \path{notes/arbitrary_fibre_hub_high_density_stability.tex} first permits
 arbitrary cut fibres and exceptional-square kernels under an adaptive
 $L^2$-energy condition, then gives a uniform exact radius with no restriction
 on either kernel whenever the complementary core is complete.  Finally,
 \path{notes/arbitrary_graphon_high_density_stability.tex} uses a
 high-degree-set decomposition to close the endpoint for every graphon.
These results change no catalogue status.

\section{Independent-set identity}

Let $\mathcal I(H)$ denote the independent vertex subsets of a finite simple
graph $H$, including the empty set.  Write $n=v(H)$ and
\[
 \Phi_H(p)=(1-p)^n\chi_H\!\left(\frac1{1-p}\right),
 \qquad
 \Gap(W)=t(H,W)-\Phi_H(p(W)).
\]

\begin{lemma}[Exact one-clique-hole density]\label{lem:independent}
For the graphon \eqref{eq:hole},
\begin{equation}
 t(H,W_s)=
 \sum_{I\in\mathcal I(H)}s^{|I|}(1-s)^{n-|I|}.
 \label{eq:independent}
\end{equation}
Hence
\begin{equation}
 \Gap(W_s)=
 \sum_{I\in\mathcal I(H)}s^{|I|}(1-s)^{n-|I|}
 -s^{2n}\chi_H(s^{-2}),
 \label{eq:gap}
\end{equation}
where the second term is interpreted by its polynomial continuation at
$s=0$.
\end{lemma}

\begin{proof}
A map $f:V(H)\to\Omega$ contributes one to the integrand of $t(H,W_s)$
unless some edge of $H$ has both endpoints mapped into $S$.  Therefore it
survives precisely when $f^{-1}(S)$ is independent.  The events
$f^{-1}(S)=I$ are disjoint and have probabilities
$s^{|I|}(1-s)^{n-|I|}$, proving \eqref{eq:independent}.  Since
$p(W_s)=1-s^2$, substitution into the definition of $\Phi_H$ gives
\eqref{eq:gap}.  Both sides are polynomials in $s$, so the formula at $s=0$
follows without a limiting interchange.
\end{proof}

This identity is valid on an arbitrary probability space.  Only the mass of
$S$ is used; no atomlessness, regularity, or finite-step approximation is
assumed.

\section{Exact Bernstein certificate}

For a polynomial $R$ of degree $d$ and a rational endpoint $b>0$, define its
Bernstein coefficients $\beta_0,\ldots,\beta_d$ by
\begin{equation}
 R(bx)=\sum_{i=0}^d\beta_i\binom di x^i(1-x)^{d-i}.
 \label{eq:bernstein}
\end{equation}
If all $\beta_i>0$, then $R(s)>0$ throughout $0\leq s\leq b$, because the
Bernstein basis functions are nonnegative and sum to one.

Exact independent-set enumeration and deletion--contraction for the
chromatic polynomial give, on every current open row,
\begin{equation}
 \Gap(W_s)=s^3(1-s)^\ell R_H(s).
 \label{eq:factor}
\end{equation}
For chromatic number three we certify $R_H>0$ on $[0,3/4]$; for chromatic
number four we use $[0,3/5]$.  These rational intervals contain the required
ones since
\[
 \frac1{\sqrt2}<\frac34,
 \qquad
 \frac1{\sqrt3}<\frac35.
\]
The following table lists every certificate.  The entries in the last
column are the complete native-degree coefficient vector
$(\beta_0,\ldots,\beta_d)$ from \eqref{eq:bernstein}; no coefficients are
omitted and no degree elevation is used.

\scriptsize
\[
\begin{array}{c|c|c|c|c|l}
\text{Atlas}&\text{graph6}&n&e&\chi(H),\ell&(\beta_0,\ldots,\beta_d)\\ \hline
43&\texttt{Dlc}&5&6&3,1&(8,101/16,121/32,133/128,79/128)\\
118&\texttt{Eht?}&6&7&3,2&(11,52/5,703/80,3863/640,889/320,937/512)\\
122&\texttt{Eld?}&6&7&3,2&(10,191/20,647/80,3511/640,761/320,809/512)\\
124&\texttt{Exd?}&6&7&3,2&(10,191/20,647/80,3511/640,761/320,809/512)\\
126&\texttt{ERUO}&6&7&3,2&(10,191/20,647/80,3511/640,761/320,809/512)\\
127&\texttt{EZEG}&6&7&3,1&(9,61/8,113/20,273/80,127/80,223/256,123/512)\\
130&\texttt{E\char`\{CW}&6&7&3,3&(8,19/2,41/4,575/64,203/64)\\
137&\texttt{EjsG}&6&8&3,2&(14,131/10,217/20,1133/160,211/80,163/128)\\
139&\texttt{Eju?}&6&8&3,2&(13,49/4,203/20,209/32,179/80,131/128)\\
145&\texttt{Exe\_}&6&8&3,1&(14,91/8,643/80,719/160,281/160,445/512,305/1024)\\
147&\texttt{EhMg}&6&8&3,1&(13,85/8,15/2,661/160,49/32,381/512,241/1024)\\
148&\texttt{EyUG}&6&8&3,1&(13,85/8,15/2,661/160,49/32,381/512,241/1024)\\
151&\texttt{EhdW}&6&8&3,1&(13,21/2,577/80,1193/320,393/320,109/128,55/256)\\
152&\texttt{EhNG}&6&8&3,1&(12,79/8,557/80,603/160,209/160,317/512,177/1024)\\
153&\texttt{Ehf\_}&6&8&3,1&(12,79/8,557/80,4851/1280,1807/1280,877/1024,165/2048)\\
157&\texttt{E\char`\~`\_}&6&9&4,2&(17,398/25,1711/125,6386/625,18889/3125,10186/3125)\\
160&\texttt{E\char`\~@g}&6&9&4,2&(16,376/25,1616/125,5986/625,17264/3125,8936/3125)\\
168&\texttt{E\char`\^MG}&6&9&3,1&(17,27/2,145/16,1423/320,75/64,371/512,215/1024)\\
169&\texttt{Exf\_}&6&9&4,1&(16,27/2,52/5,17641/2500,12846/3125,7051/3125,18208/15625)\\
171&\texttt{Ehew}&6&9&3,1&(16,51/4,341/40,1051/256,1347/1280,857/1024,113/2048)\\
174&\texttt{EtTg}&6&9&3,1&(16,101/8,659/80,4739/1280,959/1280,967/1024,71/2048)\\
178&\texttt{En\char`\}?}&6&10&4,2&(20,461/25,1921/125,13399/1250,16756/3125,7804/3125)\\
181&\texttt{E\char`\^mG}&6&10&4,1&(20,83/5,312/25,5072/625,13628/3125,6693/3125,15944/15625)\\
185&\texttt{Exv\_}&6&10&4,1&(20,33/2,1532/125,19493/2500,12738/3125,52/25,3256/3125)\\
188&\texttt{EzNG}&6&10&3,1&(20,31/2,49/5,5143/1280,499/1280,947/1024,19/2048)\\
194&\texttt{E\char`\~wW}&6&11&4,1&(24,39/2,1764/125,4269/500,2526/625,5949/3125,14352/15625)\\
196&\texttt{ER\char`\~g}&6&11&4,1&(24,39/2,1764/125,4269/500,2526/625,5949/3125,14352/15625)\\
199&\texttt{El\char`\^g}&6&11&4,1&(24,97/5,1736/125,411/50,2348/625,5756/3125,14688/15625)\\
203&\texttt{E\char`\~\char`\^G}&6&12&4,1&(28,112/5,1968/125,11201/1250,11632/3125,1041/625,2552/3125)
\end{array}
\]
\normalsize

\begin{theorem}[One-clique-hole no-hit theorem]\label{thm:no-hit}
Let $H$ be any graph in the table and let $W_s$ be \eqref{eq:hole}.  If
\[
 0\leq s\leq\frac1{\sqrt{\chi(H)-1}},
\]
equivalently if $p(W_s)$ lies in the conjectured density range, then
\[
 t(H,W_s)\geq\Phi_H(p(W_s)).
\]
Equality holds within this family only at $s=0$, or $p=1$.
\end{theorem}

\begin{proof}
The table gives $R_H(s)>0$ on a rational interval containing the required
interval.  That interval has $s<1$, so every factor in
\eqref{eq:factor} is positive for $s>0$.  Thus the gap is strictly positive
there.  At $s=0$, $W_s=1$ almost everywhere and both sides equal one; this
also handles the polynomially continued endpoint $p=1$ without division.
\end{proof}

\section{Arbitrary unions of clique blocks, with dust}

The regular-complement theorem leaves open highly nonregular singular
shapes.  A broad canonical family is a union of clique blocks, possibly with
an arbitrary dust remainder.  Let $(S_i)$ be finite or countably many
pairwise disjoint measurable sets, put $s_i=\mu(S_i)$, and define
\[
 U_{\boldsymbol s}=\sum_i\mathbf1_{S_i\times S_i}.
 \label{eq:clique-union}
\]
No assumption that $\bigcup_iS_i$ has full measure is made; $U_{\boldsymbol s}$
is zero on every cell incident to the leftover set.  Then
\[
 p(U_{\boldsymbol s})=\sum_i s_i^2.
\]
Every graph in the current open list is connected.  Hence all vertices of a
surviving homomorphism lie in one block and
\begin{equation}
 t(H,U_{\boldsymbol s})=\sum_i s_i^{v(H)}.
 \label{eq:clique-union-density}
\end{equation}

Let $m=\max_i s_i$; a positive summable sequence attains its supremum.
Since $\sum_i s_i\leq1$,
\begin{equation}
 p=\sum_i s_i^2
 \leq m^2+\left(\sum_{i:s_i\ne m}s_i\right)^2
 \leq m^2+(1-m)^2,
 \qquad p\leq m.
 \label{eq:max-mass}
\end{equation}
Throughout the required interval $p\geq1/2$, so $m\geq1/2$.  Put
$z=(1-m)/m\in[0,1]$.  The first inequality in \eqref{eq:max-mass} becomes
\[
 p\leq p_{max}(z):=\frac{1+z^2}{(1+z)^2},
 \qquad
 t(H,U_{\boldsymbol s})\geq m^n=\frac1{(1+z)^n}.
 \label{eq:clique-union-lower}
\]
For $\chi(H)=3$, the possible range is $0\leq z\leq1$.  For $\chi(H)=4$,
the condition $p_{\max}(z)\geq2/3$ is, on $[0,1]$, equivalent to
$0\leq z\leq2-\sqrt3$.

Write $p_{min}=1-1/(\chi(H)-1)$ and parametrize every possible density
between the required threshold and the upper bound in
\eqref{eq:clique-union-lower} by
\begin{equation}
 p(z,\rho)=p_{min}+\rho\bigl(p_{max}(z)-p_{min}\bigr),
 \qquad0\leq\rho\leq1.
 \label{eq:clique-density-parameter}
\end{equation}
Define the rational-coefficient polynomial
\begin{equation}
 M_H(z,\rho)=(1+z)^{2n}
 \left(\frac1{(1+z)^n}-\Phi_H(p(z,\rho))\right).
 \label{eq:clique-union-polynomial}
\end{equation}
The apparent denominators cancel; multiplying by the positive constant
$(\chi(H)-1)^n$ makes all coefficients integral.

\begin{theorem}[Clique-union-with-dust full-interval theorem]
\label{thm:clique-union}
For every current open Atlas graph $H$, every finite or countable union of
clique blocks $U_{\boldsymbol s}$, and every arbitrary leftover dust set on
which the graphon is zero, if the density lies in the conjectured range then
\[
 t(H,U_{\boldsymbol s})\geq
 \Phi_H(p(U_{\boldsymbol s})).
\]
Equality holds only for the single full-measure clique, at $p=1$.
\end{theorem}

\begin{proof}
Let $b=1$ when $\chi(H)=3$ and $b=2-\sqrt3$ when $\chi(H)=4$.  Expand
$M_H(bx,y)$ in the native bivariate Bernstein basis on $[0,1]^2$:
\begin{equation}
 M_H(bx,y)=
 \sum_{i=0}^{d_z}\sum_{j=0}^{d_\rho}
 \beta_{H,i,j}\binom{d_z}{i}\binom{d_\rho}{j}
 x^i(1-x)^{d_z-i}y^j(1-y)^{d_\rho-j}.
 \label{eq:clique-union-bernstein}
\end{equation}
Exact deletion--contraction and expansion of
\eqref{eq:clique-union-polynomial} give $20$ distinct coefficient arrays for
the $29$ rows.  Every coefficient is nonnegative in $\mathbb Q$ for
$\chi=3$ or in $\mathbb Q(\sqrt3)$ for $\chi=4$.  There is exactly one zero
coefficient in each array, namely $(i,j)=(0,d_\rho)$, corresponding to
$z=0$, $p=1$.  Positivity in the quadratic field is checked without
decimals: for $a+b\sqrt3$ with opposite signs, compare $a^2$ with $3b^2$
after recording the sign of the positive term.  The complete grouping,
degrees, and corner coefficients are
\[
\begin{array}{c|l|c|c|c}
\chi&\text{Atlas IDs}&(d_z,d_\rho)&\beta_{0,0}&\beta_{d_z,d_\rho}\\ \hline
3&43&(10,4)&8&256\\
3&118,122,124,126&(12,5)&16&1024\\
3&127&(12,5)&4&256\\
3&130&(12,5)&8&512\\
3&137,139&(12,5)&4&256\\
3&145,147,148,152&(12,5)&8&512\\
3&151&(12,5)&4&256\\
3&153&(12,5)&16&1024\\
3&168&(12,5)&8&512\\
3&171,174,188&(12,5)&16&1024\\ \hline
4&157,160&(12,5)&81&454896-262440\sqrt3\\
4&169,178&(12,5)&27&151632-87480\sqrt3\\
4&181,185&(12,5)&81&454896-262440\sqrt3\\
4&194,196&(12,5)&81&454896-262440\sqrt3\\
4&199&(12,5)&81&454896-262440\sqrt3\\
4&203&(12,5)&81&454896-262440\sqrt3
\end{array}
\]
For example,
$454896^2-3(262440)^2=306110016>0$ and
$151632^2-3(87480)^2=34012224>0$.  The verifier makes the corresponding
exact comparison for every intermediate coefficient.  Bivariate Bernstein
nonnegativity proves $M_H\geq0$ on the whole rectangle.  Moreover equality
can occur only at $x=0,y=1$: otherwise some Bernstein basis function whose
coefficient is strictly positive is nonzero.  Combine
this with \eqref{eq:clique-union-lower} and
\eqref{eq:clique-density-parameter}.  The exceptional corner is precisely
$m=p=1$, where both sides equal one.
\end{proof}

The two-clique union, whose complement is a complete bipartite cut, is a
special case.  Taking one clique block $B$ and dust $A$ gives
$W=\mathbf1_{B\times B}$; equivalently, its complement is one on every pair
incident to the exceptional hub set $A$.  Thus the theorem excludes the
canonical persistent-hub complement on the whole admissible interval.  More
generally it excludes every finite or countable clique-union-with-dust shape,
rather than only asymptotically.  These graphons are generally nonregular, so
this is independent of the regular- and maximum-degree-complement theorems.

\section{A subcritical number of deleted cliques}

The one-hole identity extends asymptotically and uniformly to a subcritical
number of deleted diagonal cliques.  Let $S_1,\ldots,S_r$ be disjoint
measurable sets of masses $s_1,\ldots,s_r$, and put
\[
 W=1-\sum_{i=1}^r\mathbf 1_{S_i\times S_i},
 \qquad \sigma_j=\sum_{i=1}^r s_i^j,
 \qquad q=1-p(W)=\sigma_2.
 \tag{4.1}\label{eq:many-holes}
\]
For $A\subseteq E(H)$, let $r_H(A)=v(H)-c(V(H),A)$ be its graphic-matroid
rank.  If $C$ runs through the nontrivial connected components of $(V(H),A)$,
define
\[
 Z(A)=\prod_C\sigma_{v(C)}.
\]
Let
\[
 P_2(H)=\sum_{x\in V(H)}\binom{d_H(x)}2,
 \qquad N_3(H)=\#\{\text{triangles of }H\},
\]
and let $N_{\geq3}(H)$ be the number of edge subsets $A$ with
$r_H(A)\geq3$.

\begin{lemma}[Exact cluster-rank identity]\label{lem:cluster-rank}
For \eqref{eq:many-holes},
\begin{equation}
 \Gap(W)=
 \bigl(P_2(H)-N_3(H)\bigr)(\sigma_3-q^2)
 +\sum_{r_H(A)\geq3}(-1)^{|A|}
 \bigl(Z(A)-q^{r_H(A)}\bigr).
 \label{eq:cluster-rank}
\end{equation}
\end{lemma}

\begin{proof}
Expand the product defining $t(H,W)$.  For a selected edge set $A$, every
nontrivial connected component must be mapped wholly into one $S_i$; the
components choose their indices independently.  Its contribution is
therefore $Z(A)$.  Whitney's expansion gives
\[
 \Phi_H(1-q)=\sum_{A\subseteq E(H)}(-1)^{|A|}q^{r_H(A)}.
\]
The rank-zero and rank-one differences vanish.  A rank-two edge set is
either two disjoint edges, two adjacent edges, or a triangle.  Disjoint
edges again give zero.  Each adjacent pair gives $\sigma_3-q^2$, while each
triangle has the opposite sign and gives $-(\sigma_3-q^2)$.  There are
$P_2(H)$ adjacent pairs and $N_3(H)$ triangles.  Collecting them proves
\eqref{eq:cluster-rank}.
\end{proof}

\begin{theorem}[Uniform subcritical clique-hole stability]
Put $c_H=P_2(H)-N_3(H)$ and suppose $c_H>0$.  For any $r\geq1$, every
graphon of the form \eqref{eq:many-holes} satisfies $\Gap(W)>0$ whenever
\begin{equation}
 0<q<\frac{c_H^2}
 {r\bigl(c_H+2N_{\geq3}(H)\bigr)^2}.
 \label{eq:many-hole-range}
\end{equation}
\end{theorem}

\begin{proof}
Because $s_i\leq\sqrt q$, for every $j\geq2$,
\[
 \sigma_j=\sum_i s_i^2s_i^{j-2}\leq q^{j/2}.
\]
If $r_H(A)\geq3$, and its nontrivial components have orders
$j_1,\ldots,j_k$, then
$\sum(j_a-1)=r_H(A)\geq3$, whence $\sum j_a\geq4$.  As $q\leq1$, this gives
\[
 0\leq Z(A)\leq q^2,
 \qquad 0\leq q^{r_H(A)}\leq q^2.
\]
Thus the absolute value of the remainder in
\eqref{eq:cluster-rank} is at most $2N_{\geq3}(H)q^2$.  The power-mean
inequality gives
\[
 \sigma_3\geq\frac{q^{3/2}}{\sqrt r}.
\]
Consequently
\[
 \Gap(W)\geq
 q^{3/2}\left(
 \frac{c_H}{\sqrt r}
 -\bigl(c_H+2N_{\geq3}(H)\bigr)\sqrt q
 \right),
\]
which is strict under \eqref{eq:many-hole-range}.
\end{proof}

For the current open rows, exact enumeration gives $c_H>0$, ranging from
$8$ to $28$, so the theorem applies to all $29$ rows.  It is uniform over
the cluster masses.  Equivalently, if
\[
 \alpha_H=\frac{c_H^2}{(c_H+2N_{\geq3}(H))^2},
\]
then the conclusion holds whenever $rq<\alpha_H$.  It therefore covers not
only every fixed number of clusters, but every family with
$r=o(1/q)$ as $q\downarrow0$.

The dependence on $r$ can in fact be removed.  This is important at the
critical Turan scale $r\asymp1/q$.

\begin{lemma}[Power-sum defect propagation]\label{lem:defect}
Let $(s_i)$ be a finite or countable sequence of nonnegative numbers with
$\sum_i s_i\leq1$.  Put
\[
 \sigma_j=\sum_i s_i^j,\qquad q=\sigma_2,\qquad
 m=\sup_i s_i,\qquad D=\sigma_3-q^2.
\]
Then $q\leq m\leq\sqrt q$, $D\geq0$, and for every integer $j\geq3$,
\begin{equation}
 0\leq \sigma_j-q^{j-1}
 \leq \binom{j-1}{2}m^{j-3}D.
 \label{eq:defect}
\end{equation}
\end{lemma}

\begin{proof}
Write $S=\sigma_1$.  Then $q\leq mS\leq m$ and $m^2\leq q$.
Interpolation between the first and $j$th power sums gives
\[
 q^{j-1}\leq S^{j-2}\sigma_j\leq\sigma_j,
\]
which proves every lower bound, including $D\geq0$.

For the upper bound, first note the exact decompositions
\[
 D=(1-S)\sigma_3+
 \sum_{i<k}s_i s_k(s_i-s_k)^2
\]
and, for $j\geq3$,
\[
 \sigma_{j+1}-q\sigma_j
 =(1-S)\sigma_{j+1}
 +\sum_{i<k}s_i s_k(s_i-s_k)
       (s_i^{j-1}-s_k^{j-1}).
\]
The quotient
$(s_i^{j-1}-s_k^{j-1})/(s_i-s_k)$ is a sum of $j-1$
monomials, each at most $m^{j-2}$.  Also
$\sigma_{j+1}\leq m^{j-2}\sigma_3$.  Hence
\begin{equation}
 0\leq\sigma_{j+1}-q\sigma_j
 \leq(j-1)m^{j-2}D.
 \label{eq:recurrence}
\end{equation}
Starting from \eqref{eq:defect} at $j=3$, use
\[
 \sigma_{j+1}-q^j
 =\sigma_{j+1}-q\sigma_j
  +q(\sigma_j-q^{j-1})
\]
and $q\leq m$.  The constants obey
$B_3=1$ and $B_{j+1}=j-1+B_j$, hence
$B_j=\binom{j-1}{2}$.  This proves \eqref{eq:defect}.
All sums are absolutely convergent; the same identities for a countable
sequence follow directly, or by truncation.
\end{proof}

For an edge subset $A$, let its nontrivial connected components have orders
$j_1,\ldots,j_k$ and rank $R=\sum_a(j_a-1)$.  Define
\[
 b(A)=\sum_{a:j_a\geq3}\binom{j_a-1}{2},
 \qquad
 \Lambda_H=\sum_{A:r_H(A)\geq3}b(A).
\]

\begin{lemma}[Uniform component-product defect]\label{lem:product-defect}
If $r_H(A)=R\geq3$, then
\begin{equation}
 0\leq Z(A)-q^R\leq b(A)mD.
 \label{eq:product-defect}
\end{equation}
\end{lemma}

\begin{proof}
Telescope the difference between
$Z(A)=\prod_a\sigma_{j_a}$ and
$q^R=\prod_a q^{j_a-1}$.  A factor with $j_a=2$ has zero difference.  For
$j_a\geq3$, Lemma~\ref{lem:defect} bounds its difference by
$\binom{j_a-1}{2}m^{j_a-3}D$.  Every unchanged factor, of either type,
is at most
\[
 \sigma_j\leq m^{j-2}q,
 \qquad q^{j-1}\leq m^{j-2}q.
\]
The corresponding telescoping term is therefore at most
\[
 \binom{j_a-1}{2}
 Dq^{k-1}m^{R-k-1}
 \leq \binom{j_a-1}{2}Dm^{R-2}
 \leq \binom{j_a-1}{2}Dm,
\]
because $q\leq m$, $R\geq3$, and $m\leq1$.  Summing the terms proves
\eqref{eq:product-defect}.
\end{proof}

\begin{theorem}[All-part-count complete-multipartite stability]
Let $H$ satisfy $c_H=P_2(H)-N_3(H)>0$.  Let $(S_i)$ be any finite or
countable family of disjoint measurable sets, and set
\[
 W=1-\sum_i\mathbf1_{S_i\times S_i},
 \qquad q=1-p(W)=\sum_i\mu(S_i)^2.
\]
If
\begin{equation}
 0<q<\left(\frac{c_H}{\Lambda_H}\right)^2,
 \label{eq:all-parts-range}
\end{equation}
then
\begin{equation}
 \Gap(W)\geq(c_H-\Lambda_H\sqrt q)D\geq0.
 \label{eq:all-parts-bound}
\end{equation}
Equality occurs only when all positive part masses are equal and cover the
whole probability space; then $q=1/k$ for an integer $k$ and $W=T_k$ up to
null sets.  Every other such graphon has strict inequality.
\end{theorem}

\begin{proof}
Combine the exact cluster-rank identity \eqref{eq:cluster-rank} with
Lemma~\ref{lem:product-defect}; signs can only decrease the rank-at-least
three remainder by $\Lambda_HmD$.  Since $m\leq\sqrt q$, this proves
\eqref{eq:all-parts-bound}, and \eqref{eq:all-parts-range} makes its first
coefficient positive.

It remains to identify $D=0$.  Equality in
$q^2\leq S\sigma_3\leq\sigma_3$ forces $S=1$ and equality in Cauchy's
inequality.  Hence all positive $s_i$ are equal.  Conversely such a
partition is the balanced complete $k$-partite graphon and every parenthesis
in \eqref{eq:cluster-rank} vanishes.
\end{proof}

For the $29$ current open rows, the exact signatures $(c_H,\Lambda_H)$ are
\[
\begin{array}{c|rrrrr}
H&43&118&122&124&126\\
(c_H,\Lambda_H)&(8,161)&(11,481)&(10,461)&(10,464)&(10,454)\\ \hline
H&127&130&137&139&145\\
(c_H,\Lambda_H)&(9,454)&(8,376)&(14,1173)&(13,1134)&(14,1211)\\ \hline
H&147&148&151&152&153\\
(c_H,\Lambda_H)&(13,1184)&(13,1207)&(13,1237)&(12,1159)&(12,1174)\\ \hline
H&157&160&168&169&171\\
(c_H,\Lambda_H)&(17,2725)&(16,2638)&(17,2893)&(16,2778)&(16,2903)\\ \hline
H&174&178&181&185&188\\
(c_H,\Lambda_H)&(16,2972)&(20,6138)&(20,6482)&(20,6598)&(20,6775)\\ \hline
H&194&196&199&203&\\
(c_H,\Lambda_H)&(24,14648)&(24,14852)&(24,14992)&(28,32316)&
\end{array}
\]
Thus the theorem covers all $29$ rows without any restriction on the number
of multipartite classes.  In particular it reaches the critical
$r\asymp1/q$ scale and includes arbitrary unbalanced part sizes.

\section{Universal rank-two complement defect}

There is also an exact statement for a completely arbitrary complement,
before any multipartite assumption is imposed.  Let $V=1-W$ and
$q=\int V=1-p(W)$.  Write
\[
 a=t(P_3,V)=\int d_V(x)^2\,dx,
 \qquad b=t(K_3,V).
\]

\begin{theorem}[Universal rank-two defect]\label{thm:rank-two}
The complete contribution of graphic ranks at most two to $\Gap(W)$ is
\begin{equation}
 \bigl(P_2(H)-N_3(H)\bigr)(a-q^2)+N_3(H)(a-b).
 \label{eq:rank-two-defect}
\end{equation}
Both defects in \eqref{eq:rank-two-defect} are nonnegative for every
graphon $V$.  If both coefficients are positive and $0<q<1$, the expression
vanishes exactly when $V$ is, up to null sets, the union of equal diagonal
clique blocks of mass $q$.  Necessarily $q=1/k$ for an integer $k$, and
$W=T_k$.
\end{theorem}

\begin{proof}
Direct inclusion--exclusion and Whitney's formula give
\[
 \Gap(W)=\sum_{A\subseteq E(H)}(-1)^{|A|}
 \bigl(t(H_A,V)-q^{r_H(A)}\bigr).
\]
Ranks zero and one vanish.  A two-edge forest contributes zero if its edges
are disjoint and contributes $a-q^2$ if they are adjacent.  There are
$P_2(H)$ adjacent pairs.  A triangle has rank two, sign minus, and contributes
$-(b-q^2)$.  This proves \eqref{eq:rank-two-defect}.

Moreover
\[
 a-q^2=\int(d_V-q)^2\geq0,
\]
while, since $0\leq V\leq1$,
\[
 a-b=\int V(x,y)V(y,z)(1-V(x,z))\,dx\,dy\,dz\geq0.
\]

Suppose both vanish.  Then $d_V=q$ almost everywhere, and for almost every
triple
\begin{equation}
 V(x,y)V(y,z)>0\quad\Longrightarrow\quad V(x,z)=1.
 \label{eq:transitivity}
\end{equation}
For a typical $x$, let $N_x=\{z:V(x,z)>0\}$.  If $y\in N_x$ is typical,
apply \eqref{eq:transitivity} to $(z,x,y)$: for almost every $z\in N_x$,
$V(z,y)=1$.  Hence
\[
 q=d_V(x)\leq\mu(N_x)\leq d_V(y)=q.
\]
All inequalities are equalities.  Thus $V(x,\cdot)=1$ almost everywhere on
$N_x$, and for almost every $y\in N_x$ one has $N_y=N_x$ modulo null sets.
The positive-support relation therefore partitions the probability space
into clique blocks, each of measure $q$.  Constant degree leaves no positive
measure outside the blocks, so their number is $1/q$ and they are balanced.
The converse is immediate.
\end{proof}

For every current open row, both coefficients in
\eqref{eq:rank-two-defect} are strict:
\[
 8\leq P_2(H)-N_3(H)\leq28,
 \qquad 1\leq N_3(H)\leq9.
\]
At the level of this rank-two theorem alone, an arbitrary high-density
negative sequence would have to make both universal defects small enough for
graphic-rank-at-least-three terms to compete.  The exact common zero set is
precisely the balanced Turan family, and the preceding all-part-count theorem
controls every exact complete-multipartite complement near that set.

The independent note
\path{notes/stochastic_block_high_density_stability.tex} goes further: its
operator-spectral cycle-defect estimate controls every arbitrary measurable
complement of constant degree $q$, proving strict positivity uniformly for
sufficiently small $q$.  No equal-cell or finite-step assumption is needed.
Symmetric doubly stochastic equal-cell matrices are a corollary; on their
canonical segment from $I_r$ to $J_r/r$, an exact bivariate Bernstein theorem
removes the high-density restriction and works for every
$r\geq\chi(H)-1$.

\section{Interpretation, assumptions, and limitations}

The leading factor $s^3=(1-p)^{3/2}$ is consistent with the triangle term in
the cycle-rank expansion, but the proof above is exact on the full admissible
interval rather than merely asymptotic.  It closes the simplest singular
escape left by bounded normalized complements: one cluster carries all the
missing edges at amplitude one.

The exact full-interval calculation covers one deleted cluster.  The final
multi-cluster theorem covers every finite or countable complete-multipartite
partition, including the critical scale $r\asymp1/q$, but only in its
explicit high-density interval.  The independent bivariate theorem covers
the complementary family of every finite or countable union of clique blocks
on the full interval.  These statements alone do not cover arbitrary shapes
outside both cluster families.  The later note
\path{notes/arbitrary_graphon_high_density_stability.tex} avoids a structural
classification: it splits at the high-degree set and combines the
complete-core and maximum-degree margins.  This closes a uniform endpoint
interval for every arbitrary graphon.  Balanced complete multipartite
graphons remain equality cases, as required.  The results are endpoint
no-hit tests and do not prove any Atlas row on its full density interval.

The only inputs are the exact identity of Lemma~\ref{lem:independent}, the
maximum-mass bound \eqref{eq:max-mass}, exact independent-set enumeration,
exact deletion--contraction for $\chi_H$, and the elementary univariate and
bivariate Bernstein positivity criteria.  There is no numerical root claim,
optimization, graphon approximation, conditioning, or unsafe denominator.
The graph6 codes identify every graph unambiguously; the verifier also checks
its order, size, chromatic number, factor multiplicities, and every rational
coefficient displayed above.

Run
\begin{verbatim}
python codes/principled_negative_tests.py
\end{verbatim}
from the repository root.  The line headed ``singular
one-clique-complement full-interval certificates'' reconstructs all $29$
factorizations and the $25$ distinct Bernstein vectors.  The line headed
``nonregular clique-union'' reconstructs the
$20$ bivariate coefficient classes over $\mathbb Q$ and
$\mathbb Q(\sqrt3)$.  The same run reconstructs all $29$
subcritical-cluster, all-part-count, and regular-complement signatures using
exact arithmetic.

\end{document}
